\section{问题二：历史误差场景下的滚动日前调度}
\subsection{预测与有效残差样本}
日前决策不能使用当天实际负载和光伏。负载预测采用最近4次同星期几的均值，并叠加最近5天预测残差的均值；光伏预测采用最近7天均值和最近28天残差漂移。将两者转换为电量后相减，得到净负载点预测$\widehat N_{d,t}$。所有残差均以该历史日当时可得的预测计算，而不是用事后拟合值替代。

场景集合取最近56天内的有效历史日$\mathcal H_d$。对$r_s\in\mathcal H_d$，
\begin{equation}
N_{s,t}=\widehat N_{d,t}+N_{r_s,t}-\widehat N_{r_s,t},\qquad
\omega_s=1/|\mathcal H_d|.\label{eq:residual}
\end{equation}
前7天缺少同星期历史，负载预测使用典型日回退，因此最终方案从1月8日起纳入残差。早期回退误差本身是合法的因果预测误差，排除它们是为了使场景样本与后续预测器口径更一致。2月1日可用24条残差，随后逐渐增至56条。主模型原有预测处理保持不变，不把其他预测器改动混入该对照。

\subsection{两阶段缺口补救与48小时前瞻}
普通日采用288个时段的窗口。当天购电和充放电是场景间共享的决策，紧急购电$e_{s,t}$作为场景相关的补救变量；次日使用点预测近似。优化目标为
\begin{equation}
\min\ J_2^{\mathrm{look}}=
\sum_{k=1}^{288}p_k g_k+
5\sum_s\omega_s\sum_{t=1}^{144}p_t e_{s,t}.\label{eq:q2}
\end{equation}
其中$p_k$为重复至两日的典型日电价。需求约束为
\begin{align}
g_k+b_k-c_k&\ge\widehat N_k,\qquad g_k\ge0,\\
g_t+b_t-c_t+e_{s,t}&\ge N_{s,t},\qquad e_{s,t}\ge0.
\end{align}
储能满足共同约束，普通日窗口末端回到本次窗口起始储电量。次日前瞻用于估计库存价值，不代表次日计划已成交；当天执行完毕后重新计算后续计划。12月31日仅保留当天，并固定年末电量6000 kWh。

在固定充放电、单时段正电价且没有其他约束时，$pg+5p\E[(N-g-b+c)_+]$的内部最优条件对应80\%分位数。这有助于解释高价缺口补救为何推动计划购电增加，但完整模型还含储能耦合、点预测供给约束和终端条件，不能逐段直接套用该分位数替代求解。

\subsection{执行流程与物理可行性}
每天首先更新预测与历史场景，求解稀疏LP，然后只执行当天144段，传递实际日末储电量。实际缺口与费用分别为
\begin{equation}
e_t=\max(N_t-g_t-b_t+c_t,0),\qquad
J_{2,d}=\sum_t p_tg_t+5\sum_t p_te_t.\label{eq:q2settle}
\end{equation}
计划购电费与未用电量无关，故降低紧急购电量并非唯一优化目标。

对于允许弃置多余供给的随机模型，设$\rho=\eta_c\eta_d<1$。如出现同时充放电，取$a=\min(c_t,b_t/\rho)$并令$c'_t=c_t-a$、$b'_t=b_t-\rho a$。该变换保持式\eqref{eq:storage}不变，净供给增加$(1-\rho)a$，功率不增，且至少一个充放电量变为零。允许弃置时可保持可行，紧急购电不会增加。主程序采用这一等库存处理并核验执行结果。该推理依赖相应弃置条件，不能无条件推广到其他设备模型。

\subsection{结果与窗口规则比较}
\begin{table}[htbp]\centering\small
\caption{第二问两种残差口径下的费用（元）}
\begin{tabular}{lrr}\toprule
费用项 & 含回退残差 & 有效残差主方案\\\midrule
计划费 & 13508190.55 & 13462061.30\\
紧急费 & 893806.43 & 927073.78\\
总费用 & 14401996.98 & 14389135.08\\
\bottomrule\end{tabular}\end{table}

主方案总费用为14389135.08元。相对含回退残差的56天等权方案，计划费减少46129.24元，紧急费增加33267.35元，总费用净降12861.89元。表中各金额分别由未舍入值计算后保留两位小数，个别差额与已舍入单元格相减可差0.01元。

\fig{03_q2_residuals}
图\ref{fig:03_q2_residuals}采用相同规则在两种残差口径下的配对结果。在有效残差条件下，56天等权费用低于半衰期28天与四候选在线选择；据此保留较简单的等权方案。该比较是同年度回顾性分析，不能证明56天在其他年份仍最优。

从时间分布看，2月节省15521.02元，3月增加2659.13元，最后一个日费用有差异的日期为3月4日。3月5日至12月31日逐日费用一致，说明改动主要作用于启动阶段的场景构成，而非全年持续提高适应性。指定日期结果见附录\ref{sec:outputs}。
\FloatBarrier
